
\chapter{ניתוח מעמיק של שגיאות ומקרי קצה}

\section{התנהגות בתדרים גבוהים}

אחד מאתגרי חילוץ גלי סינוס הוא הדיוק בתדרים גבוהים (\textenglish{$f > 3000$ Hz}), שם הספקטרוגרמה דחוסה ואותות שכנים מתנגשים. ניתחנו \textenglish{500} מקרי כשל של ה-\textenglish{BiLSTM} ומצאנו שלשה תבניות עיקריות:

\begin{itemize}
\item \textbf{ערבוב תדרים (\textenglish{Frequency Confusion})}: \textenglish{68\%} מהשגיאות אירעו כאשר \(\Delta f < 50\) הרץ בין שני מרכיבים. הפתרון המוצע הוא הגדלת רזולוציית ה-\textenglish{STFT} (חלון גדול יותר) בשיטת \textenglish{multi-resolution STFT}.
\item \textbf{אובדן פאזה (\textenglish{Phase Drift})}: בתנאי \textenglish{SNR} נמוך מ-\textenglish{$-3$ dB}, שגיאת הפאזה המשוחזרת גדלה מ-\textenglish{5°} ל-\textenglish{23°} בממוצע.
\item \textbf{משרעת נמוכה (\textenglish{Weak Component})}: כאשר יחס המשרעות \(A_1 / A_2 > 10\), המרכיב החלש לעתים "נעלם" מהפלט.
\end{itemize}

\section{השוואה עם גישות ספקטרליות קלאסיות}

בהשוואה לאלגוריתם \textenglish{MUSIC} (\textenglish{MUltiple SIgnal Classification}) ולשיטת \textenglish{ESPRIT}, ה-\textenglish{BiLSTM} הוכיח עדיפות ב-\textenglish{SNR} נמוך (\textenglish{$< 0$ dB}) אך נחיתות ב-\textenglish{SNR} גבוה (\textenglish{$> 20$ dB}) מבחינת דיוק תדר (\textenglish{$\pm 0.5$ Hz לעומת $\pm 0.1$ Hz}). הסיבה: אלגוריתמים קלאסיים מבוססים על עיקרון \textenglish{maximum likelihood} שמנצל הנחות סטטיסטיות חזקות; הרשת הנוירונית גמישה יותר אך פחות מדויקת בתנאים אידיאליים \cite{williamson2016complex}.

\section{ניסויי \textenglish{Ablation}}

\begin{table}[H]
\centering
\caption{ניסויי \textenglish{Ablation}: השפעת כל רכיב על \textenglish{SI-SNR}}
\label{tab:ablation}
\begin{english}
\begin{tabular}{lc}
\toprule
\textbf{Configuration} & \textbf{SI-SNR (dB)} \\
\midrule
Full BiLSTM + Attention + SI-SNR loss   & \textbf{11.3} \\
BiLSTM + MSE loss (no SI-SNR)           & 9.8  \\
Uni-LSTM + Attention                    & 9.1  \\
BiLSTM, no Attention                    & 10.4 \\
BiLSTM, no Curriculum                  & 10.6 \\
\bottomrule
\end{tabular}
\end{english}
\end{table}
